\section{Introduction}\label{sec:intro}

Precision oncology depends on identifying tumor-specific molecular alterations (e.g., point mutations, copy-number alterations, and gene fusions) that can be used for diagnosis, prognosis, treatment selection, and trial matching. 
These alterations are not merely associative markers: they are the molecular mechanisms driving tumor growth and progression, making them direct targets for therapy~\cite{sanchez-vegaOncogenicSignalingPathways2018, drilonEfficacyLarotrectinibTRK2018, marabelleEfficacyPembrolizumabPatients2020, mateoDeliveringPrecisionOncology2022}.

In practice, these alterations are diagnosed using molecular pathology assays that range from targeted single-gene tests to broad next-generation sequencing (NGS) panels, whole-exome sequencing, or whole-genome sequencing. 
Although the cost per sequencing run has decreased over time, routine deployment of genome-scale testing still requires substantial investment in bioinformatics infrastructure (compute, storage, pipelines) and specialized personnel to run and maintain these systems.
In their tertiary-hospital implementation study, Kim et al.~\cite{kimClinicalImplementationNextgeneration2025} explicitly highlight both the financial investment required and the bioinformatics specialists and server engineers needed to support NGS in routine care. 
More broadly, the rapid clinical uptake of NGS has created persistent downstream bottlenecks in data processing, quality control, storage, and, most importantly, clinical interpretation of complex variant findings~\cite{karakoyunChallengesClinicalInterpretation2023, remonPrecisionOncologySeparating2018}. 
These deployment barriers (cost and specialized expertise) explain the growing interest in lower-cost triage signals that prioritize which patients warrant confirmatory molecular testing, particularly in settings where universal sequencing is not feasible.

One promising candidate signal is tissue morphology: across several well-studied cancers, driver mutations leave reproducible imprints on H\&E morphology, though the strength, mechanism, and discoverability of these imprints vary by gene, tissue, and cellular context. 
Some associations are mechanistically obligate and pathologist-recognizable 
(CDH1 $\to$ discohesive/lobular~\cite{frixenEcadherinmediatedCellcellAdhesion1991, derksenSomaticInactivationEcadherin2006}; 
VHL $\to$ clear cell~\cite{qiuHIF2aDependentLipidStorage2015, duHIFDrivesLipid2017, chenGYS1InducesGlycogen2020}; 
PML-RARA $\to$ hypergranular promyelocytes~\cite{grignaniFusionProteinsRetinoic1998, linRoleHistoneDeacetylase1998, swerdlowWHOClassificationTumours, sehgalAuerRodsFaggot2023}; 
MYC $\to$ starry-sky Burkitt~\cite{evanInductionApoptosisFibroblasts1992, eischenDisruptionARFMdm21999}). 
For most other recurrent driver mutations (KRAS, EGFR, PIK3CA, TP53, ATM, etc.), the morphological correlates are more subtle and confounded by the cell-of-origin and the subtype context in which the mutation occurs.
These genotype-to-morphology relationships open a complementary triage axis: 
routine H\&E morphology can act as a low-cost triage layer that prioritizes which patients warrant confirmatory molecular testing, provided the imprints, including the subtle and confounded ones, can be extracted reliably and explained in pathologist-interpretable terms.

The last decade has seen a rapid expansion of digital pathology, driven by improvements in whole-slide imaging (WSI) scanners, storage/compute infrastructure, and clinical validation efforts~\cite{vanderlaakDeepLearningHistopathology2021, evansUSFoodDrug2018, mukhopadhyayWholeSlideImaging2018}.
These developments have made computational pathology a natural domain for deep neural networks (DNNs), because DNNs can learn predictive features directly from pixel data at multiple spatial scales~\cite{dimitriouDeepLearningWhole2019}. 
Early work established strong proof-of-concept that clinically relevant molecular alterations are, to some extent, predictable from routine hematoxylin-and-eosin (H\&E) morphology using convolutional neural networks (CNNs) trained on tiled WSIs. 
Coudray et al.~\cite{coudrayClassificationMutationPrediction2018} train a CNN on The Cancer Genome Atlas (TCGA) WSIs and reported that multiple recurrent lung adenocarcinoma gene mutations (including EGFR, KRAS, STK11, and TP53) could be predicted from H\&E with meaningful discrimination performance.
Kather et al.~\cite{katherPancancerImagebasedDetection2020} extend this approach to 14 cancer types in over 5,000 patients, with mean cross-validated AUROCs for the top driver mutations clustering in the 0.6--0.8 range across lung, colorectal, breast, and gastric cancers.
Fu et al.~\cite{fuPancancerComputationalHistopathology2020} show that whole-genome doubling has a universal morphological signature across cancer types. 

Translation to routine clinical workflows, however, has been limited, both because external-cohort generalization remains uneven~\cite{niehuesGeneralizableBiomarkerPrediction2023, arslanSystematicPancancerStudy2024} and because established biomarkers are already covered by mature IHC and ISH assays. 
A separate concern, central to the present work, is that these models offer little mechanistic insight into the genotype-morphology relationships they exploit.
Subsequent work using multiple instance learning (MIL) and attention-based aggregation~\cite{ilseAttentionbasedDeepMultiple2018, luDataefficientWeaklySupervised2021, campanellaClinicalgradeComputationalPathology2019} improve discrimination performance and produced post-hoc attention maps that purport to localize the morphological evidence behind each prediction, although the fidelity of these maps to underlying biology is contested.

A recurring observation across theses efforts is that molecular subtypes are substantially more predictable from H\&E than the individual driver mutations they contain.
Strong subtype-from-H\&E results have now been reproduced for breast intrinsic subtypes~\cite{coutureImageAnalysisDeep2018, naikDeepLearningenabledBreast2020}, colorectal consensus molecular subtypes (CMS)~\cite{bilalDevelopmentValidationWeakly2021, wagnerTransformerbasedBiomarkerPrediction2023}, endometrial ProMisE classes~\cite{volinsky-fremondPredictionRecurrenceRisk2024}, and bladder muscle-invasive subtypes~\cite{woerlDeepLearningPredicts2020}.
Kather et al.~\cite{katherPancancerImagebasedDetection2020} make the same point in their own data, finding that integrated features such as MSI, homologous recombination deficiency, and expression-based subtypes were markedly more predictable than the individual mutations comprising them.
This asymmetry has a straightforward interpretation: each subtype integrates many co-occurring molecular features that push morphology in correlated directions, producing a stronger and more reproducible histologic signal than any single driver. 
A methodological consequence is that some apparent single-mutation predictions are partly subtype prediction in disguise: a model nominally calling TP53 or PIK3CA in breast cancer recovers, in part, the basal-like or luminal-A context in which those mutations are enriched.

A second, more recent shift is the rise of pathology foundation models: 
large self-supervised vision encoders pre-trained on millions of whole-slide 
images across institutions and tissues. 
UNI~\cite{chenGeneralpurposeFoundationModel2024} (100M tiles, 100K slides, 20 tissue types), Prov-GigaPath~\cite{xuWholeslideFoundationModel2024} (1.3B tiles from 171K real-world non-TCGA slides), Virchow~\cite{vorontsovVirchowMillionSlideDigital2024} (a 632M-parameter ViT trained on 1.5M Memorial Sloan Kettering slides), and the open-weights Phikon~\cite{filiotScalingSelfSupervisedLearning2024} have collectively reset the floor for what counts as a strong off-the-shelf representation in computational pathology. 
The vision-language variant CONCH~\cite{luVisuallanguageFoundationModel2024} reports zero-shot subtype classification accuracies above 90\% on NSCLC, RCC, and breast cancer. 
These models substantially mitigate the representation-learning and external-cohort generalization problems that limited the previous CNN/MIL generation. 
They do not, however, resolve the explainability problem: hierarchical attention, joint vision-language pre-training, and patch-level pooling all complicate post-hoc attribution to specific morphological structures, and large-scale pre-training does not by itself guarantee that the features driving a prediction correspond to the biology a pathologist would recognize.


% TODO: contributions (not now, need to wait for the results)
